\section{A latent-jobs model of household labour supply}\label{a-latent-jobs-model-of-household-labour-supply}

A job is a package \(j\): an employment state \(e\in\{0,1\}\) and, when employed, an occupation \(k\in\{1,\dots,4\}\), weekly hours \(h\in[5,70]\) and an hourly wage \(w\in[2,590]\). The base measure \(\nu\) places counting mass on the single non-employment point and, on the employed part, counting measure over occupations with Lebesgue measure over hours and wages. A household chooses one package; a couple chooses one pair of packages under a shared budget.

\subsection{Preferences}\label{preferences}

Let \(C_i(j)\) be priced monthly disposable consumption at \(j\) and \(\ell_i(j)\) normalized leisure, \((80-h)/10\) hours per week, floored at one hour. For a single adult of sex group \(g\), deterministic utility is

\[
u_{ij}=\underbrace{\beta_{\ell}^{g}(\mathbf x_i)\;
\mathcal B\!\left(\tilde\ell_{ij};\theta_{\ell}^{g}\right)}_{L_{ij}}
\;+\;\beta_c\log\tilde c_{ij},
\qquad
\beta_{\ell}^{g}(\mathbf x_i)=\beta_{\ell 0}^{g}+\beta_{\ell a}^{g}a_i
+\beta_{\ell a^{2}}^{g}a_i^{2}
+\mathbb 1\{g=\text{women}\}\,\beta_{\ell k}^{g}k_i,
\]

where \(\mathcal B(z;\theta)=(z^{\theta}-1)/\theta\) with \(\mathcal B(z;0)=\log z\), \(\tilde\ell_{ij}\) is normalized leisure, \(\tilde c_{ij}=C_{ij}/\lambda_c\), \(a_i\) is centred age in decades and \(k_i\) is the child count. For a couple the index is additive across spouses with no direct cross-leisure term,

\[
u_{ij}=\beta_{\ell}^{m}(\mathbf x_i)\,\mathcal B(\tilde\ell^{m}_{ij};\theta_{\ell}^{m})
      +\beta_{\ell}^{f}(\mathbf x_i)\,\mathcal B(\tilde\ell^{f}_{ij};\theta_{\ell}^{f})
      +\beta_c\log\tilde c_{ij},
\]

with consumption the tax-unit sum over all members. We write \(L_i(j)\) throughout for the complete non-consumption part of the index, which for a couple contains both leisure terms.

Three features of this specification are decisions and are worth naming. The consumption curvature is exactly zero, so the consumption term is \(\beta_c\log(C/\lambda_c)\) and the marginal utility of consumption is \(\beta_c/C\), independent of the normalizer. The consumption weight \(\beta_c\) is \textbf{estimated}, not fixed at one. And the random-utility shock is i.i.d. type-I extreme value with scale one, so utility is measured in the natural unit of that scale. Fixing the shock scale is a normalization; additionally fixing \(\beta_c\) would be a substantive restriction relative to that normalization, and estimating \(\beta_c\) removes it. Neither step identifies an absolute cardinal utility scale, and fixing the consumption functional form is not by itself what identifies the shock scale.

The normalizer \(\lambda_c\) is a units convention. Under exact log consumption it enters the index as the alternative-invariant constant \(-\beta_c\log\lambda_c\), so it cancels from every choice probability and, as shown below, exactly from the welfare measure.

\subsection{Opportunities}\label{opportunities}

The opportunity density gives the relative intensity with which packages are available to household \(i\), before any choice is made. It is a product of \textbf{four} factors, switched off on non-employment by the employment indicator \(E_{ij}=\mathbb 1\{h_{ij}>0\}\):

\[
g_{ij}=g^{E}_{ij}\cdot\left(g^{H}_{ij}\cdot g^{\mathrm{Occ}}_{ij}\cdot g^{W}_{ij}\right)^{E_{ij}} .
\]

\textbf{Access}, \(g^{E}_{ij}=\exp\{\beta_E+\beta_s s_i+\sum_{r=2}^{8}\beta_r\,\mathrm{reg}_{ir}+\beta_u u_i+\beta_m m_i\}\), carries the local group unemployment measure \(s_i\), the region indicators and urbanisation, and is the factor inside which local-market access lives. \textbf{Hours}, \(g^{H}_{ij}=\exp\{\sum_b \beta_b\mathbb 1\{h_{ij}\in b\}\}\), elevates five bands over a residual reference set of total width 26.5 hours per week, so every band coefficient is read against that residual and none of the five is normalized to zero; the narrow full-time band \([33.5,36.5)\) is a density elevation over an interval, not an atom at thirty-five hours. \textbf{Occupation}, \(g^{\mathrm{Occ}}_{ij}=\exp\{\beta^{\mathrm{occ}}_{k,g}\}\) with \(\beta^{\mathrm{occ}}_{1,g}\equiv 0\). \textbf{Wage offer}, \(g^{W}_{ij}\), is log-normal with location \(\mu_i=\beta_{w0}+\beta_{wL}L_i+\beta_{wH}H_i+\beta_{wx}x_i+\beta_{wx^{2}}x_i^{2}+\delta_{\mathrm{occ}}\) and dispersion \(\sigma\), truncated to \([2,590]\) euros per hour and renormalized on that support.

Two properties of the normalization matter later. Write \(\widehat g_{ij}=g_{ij}/Z_i\) for the density normalized to unit mass on the whole package space. The wage factor integrates to one \emph{conditional on employment and occupation}, because the truncated density is renormalized on its own support; and the non-consumption index \(L_{ij}\) has no wage argument. Together these give the reference the pay-neutrality property established in Section 4.

The four blocks are what the decomposition later separates. Employment and occupation access are one mechanism; the distribution of wages conditional on an occupation is another; neither is the same as non-labour resources or household needs.

\subsection{Identification, in words}\label{identification-in-words}

Choices alone do not separate a taste for leisure from a scarcity of jobs at that number of hours. Three kinds of restriction do the work. Preferences are smooth in hours through the leisure index, while the opportunity density is a step function over bands, so a spike in observed hours at a band is read as availability rather than as a kink in tastes. Excluded shifters enter availability and not preferences: the local group unemployment measure and the regional and urbanisation indicators shift the employment index and appear nowhere in utility. And the wage-offer distribution is specified as independent of offered hours conditional on occupation, which separates the wage location from the hours density. These are the restrictions of \citet{capeau2016} and \citet{dagsvikjia2016}, maintained here rather than tested. The regional variation supports the access block conditional on those restrictions; it does not establish separate identification on its own, and it is not a causal design.

\subsection{Estimation}\label{estimation}

The choice set is a continuum, so the likelihood is evaluated over sampled alternatives. For each household we draw \(R=100\) alternatives from a proposal density \(q_{ij}\) and place the observed package in the set as well, giving a set \(\mathcal C_i\) with \(|\mathcal C_i|=101\) rows. Write the sampled-set index

\[
V_{ij}=u_{ij}+\log g_{ij}-\log q_{ij},
\qquad
q_{ij}=q^{E}_{ij}\left(q^{H}_{ij}q^{W}_{ij}q^{\mathrm{Occ}}_{ij}\right)^{E_{ij}} .
\]

The contribution of household \(i\) is then the conditional probability of its observed package \(y\) within that set,

\[
\Pr\!\left(y\mid\mathcal C_i\right)=
\frac{n_{y}\exp V_{iy}}{\sum_{s\in\mathcal C_i}\exp V_{is}} .
\]

The denominator runs over slots rather than distinct packages, and a package drawn more than once keeps its multiplicity, \(n_y\); the observed package carries the same proposal correction as any other row. The proposal density is a computational device and never an economic object: it is not an offer, an availability or an opportunity. The proposal is fitted out of fold, so a household's own outcome never enters the proposal used to score it. For couples the proposal draws the joint participation regime first and then the two spouse packages conditional on it.

Standard errors are cluster-robust on the household. Optimization uses five starts under two polishing contracts, ten terminal paths in all; we report the spread of the criterion across those paths and the eigenvalues of the exact Hessian at the selected optimum. Those diagnostics support a stable local solution found from the starts tested. They are not a proof of global uniqueness, and we do not claim one.

Two features of the observation rule remain genuinely open and are stated here rather than in a footnote. The sample screens on the observed hours and wage of employed deciders, so the estimation sample is selected on an outcome of the process being modelled. The conditional likelihood above is the correct object for a sampled choice set; it is not automatically the correct object for an outcome-selected sample, and we have not established which correction, if any, the screen requires. Separately, the derivation of the sampled-set probability conditions on a labelled collection of slots; small observed cross-coordinate correlations and a binomial-looking multiplicity distribution are consistent with independent draws but do not by themselves establish the sampling law. Both are recorded as unresolved in Section 6.

